%%
%% This is file `main.tex' based on `sample-sigconf.tex' (q.v. for spurce of that,
%%
%% IMPORTANT NOTICE:
%% 
%% For the copyright see the original source file `sample-sigconf.tex'
%% in the `Sample' folder.
%%
%% For distribution of the original source see the terms
%% for copying and modification in the file samples.dtx.
%% 
%% This generated file may be distributed as long as the
%% original source files, as listed above, are part of the
%% same distribution. (The sources need not necessarily be
%% in the same archive or directory.)
%%
%% Commands for TeXCount
%TC:macro \cite [option:text,text]
%TC:macro \citep [option:text,text]
%TC:macro \citet [option:text,text]
%TC:envir table 0 1
%TC:envir table* 0 1
%TC:envir tabular [ignore] word
%TC:envir displaymath 0 word
%TC:envir math 0 word
%TC:envir comment 0 0
%%
%%
%% The first command in your LaTeX source must be the \documentclass command.

%% NOTE that a single column version is required for 
%% submission and peer review. This can be done by changing
%% the \doucmentclass[...]{acmart} in this template to 
%%\documentclass[manuscript,review,anonymous]{acmart}
%% This version is used for drafting and final submission
\documentclass[sigconf]{acmart}

%% 
%% To ensure 100% compatibility, please check the white list of
%% approved LaTeX packages to be used with the Master Article Template at
%% https://www.acm.org/publications/taps/whitelist-of-latex-packages 
%% before creating your document. The white list page provides 
%% information on how to submit additional LaTeX packages for 
%% review and adoption.
%% Fonts used in the template cannot be substituted; margin 
%% adjustments are not allowed.

%%
%% \BibTeX command to typeset BibTeX logo in the docs
\AtBeginDocument{%
  \providecommand\BibTeX{{%
    \normalfont B\kern-0.5em{\scshape i\kern-0.25em b}\kern-0.8em\TeX}}}

%% Rights management information.  This information is sent to you
%% when you complete the rights form.  These commands have SAMPLE
%% values in them; it is your responsibility as an author to replace
%% the commands and values with those provided to you when you
%% complete the rights form.
\setcopyright{acmlicensed}
\copyrightyear{2024}
\acmYear{2024}
\acmDOI{XXXXXXX.XXXXXXX}

%% These commands are for a PROCEEDINGS abstract or paper.
\acmConference[Conference acronym 'XX]{Make sure to enter the correct
  conference title from your rights confirmation email}{June 03--05,
  2018}{Woodstock, NY}
%
%  Uncomment \acmBooktitle if th title of the proceedings is different
%  from ``Proceedings of ...''!
%
%\acmBooktitle{Woodstock '18: ACM Symposium on Neural Gaze Detection,
%  June 03--05, 2018, Woodstock, NY} 
\acmISBN{978-1-4503-XXXX-X/18/06}


%%
%% Submission ID.
%% Use this when submitting an article to a sponsored event. You'll
%% receive a unique submission ID from the organizers
%% of the event, and this ID should be used as the parameter to this command.
%%\acmSubmissionID{123-A56-BU3}

%%
%% For managing citations, it is recommended to use bibliography
%% files in BibTeX format.
%%
%% You can then either use BibTeX with the ACM-Reference-Format style,
%% or BibLaTeX with the acmnumeric or acmauthoryear sytles, that include
%% support for advanced citation of software artefact from the
%% biblatex-software package, also separately available on CTAN.
%%
%% Look at the sample-*-biblatex.tex files for templates showcasing
%% the biblatex styles.
%%

%%
%% The majority of ACM publications use numbered citations and
%% references.  The command \citestyle{authoryear} switches to the
%% "author year" style.
%%
%% If you are preparing content for an event
%% sponsored by ACM SIGGRAPH, you must use the "author year" style of
%% citations and references.
%% Uncommenting
%% the next command will enable that style.
%%\citestyle{acmauthoryear}

%%
%% end of the preamble, start of the body of the document source.

%% % Location of your graphics files for figures, here a sub-folder to the main project folder
\graphicspath{{./images/}}

\usepackage{listings}
\usepackage{booktabs}

\begin{document}

%%
%% The "title" command has an optional parameter,
%% allowing the author to define a "short title" to be used in page headers.
\title{Project report}

%%
%% The "author" command and its associated commands are used to define
%% the authors and their affiliations.
%% Of note is the shared affiliation of the first two authors, and the
%% "authornote" and "authornotemark" commands
%% used to denote shared contribution to the research.

\author{Daniel Granstr{\"o}m}
\affiliation{%
\institution{Aalto University}
\streetaddress{1 Th{\o}rv{\"a}ld Circle}
\city{Espoo}
\country{Finland}}
\email{daniel.granstrom@aalto.fi}


\author{Joonatan Korpela}
\affiliation{%
  \institution{Aalto University}
  \city{Espoo}
  \country{Finland}
}
\email{joonatan.korpela@aalto.fi}

\author{Mikael Siidorow}
\affiliation{%
  \institution{Aalto University}
  \city{Espoo}
  \country{Finland}
}
\email{mikael.siidorow@aalto.fi}

%%
%% By default, the full list of authors will be used in the page
%% headers. Often, this list is too long, and will overlap
%% other information printed in the page headers. This command allows
%% the author to define a more concise list
%% of authors' names for this purpose.
\renewcommand{\shortauthors}{Trovato and Tobin, et al.}

%%
%% The abstract is a short summary of the work to be presented in the
%% article.
\begin{abstract}
  A clear and well-documented \LaTeX\ document is presented as an
  article formatted for publication by ACM in a conference proceedings
  or journal publication. Based on the ``acmart'' document class, this
  article presents and explains many of the common variations, as well
  as many of the formatting elements an author may use in the
  preparation of the documentation of their work.
\end{abstract}

%%
%% The code below is generated by the tool at: http://dl.acm.org/ccs.cfm
%% Please copy and paste the code instead of the example below.
%%
\begin{CCSXML}
  <ccs2012>
  <concept>
  <concept_id>00000000.0000000.0000000</concept_id>
  <concept_desc>Do Not Use This Code, Generate the Correct Terms for Your Paper</concept_desc>
  <concept_significance>500</concept_significance>
  </concept>
  <concept>
  <concept_id>00000000.00000000.00000000</concept_id>
  <concept_desc>Do Not Use This Code, Generate the Correct Terms for Your Paper</concept_desc>
  <concept_significance>300</concept_significance>
  </concept>
  <concept>
  <concept_id>00000000.00000000.00000000</concept_id>
  <concept_desc>Do Not Use This Code, Generate the Correct Terms for Your Paper</concept_desc>
  <concept_significance>100</concept_significance>
  </concept>
  <concept>
  <concept_id>00000000.00000000.00000000</concept_id>
  <concept_desc>Do Not Use This Code, Generate the Correct Terms for Your Paper</concept_desc>
  <concept_significance>100</concept_significance>
  </concept>
  </ccs2012>
\end{CCSXML}

\ccsdesc[500]{Do Not Use This Code~Generate the Correct Terms for Your Paper}
\ccsdesc[300]{Do Not Use This Code~Generate the Correct Terms for Your Paper}
\ccsdesc{Do Not Use This Code~Generate the Correct Terms for Your Paper}
\ccsdesc[100]{Do Not Use This Code~Generate the Correct Terms for Your Paper}

%%
%% Keywords. The author(s) should pick words that accurately describe
%% the work being presented. Separate the keywords with commas.
\keywords{Do, Not, Us, This, Code, Put, the, Correct, Terms, for,
  Your, Paper}

%% The following are not a requirement, delete if not using
\received{20 February 2024}  %% inital submission date
\received[revised]{12 March 2024} %% interim new draft
\received[accepted]{5 June 2024}  %% publication version

%%
%% This command processes the author and affiliation and title
%% information and builds the first part of the formatted document.
\maketitle

\section{Introduction}
Papers do not have any fixed sections names. Three levels of heading are supported in ACM papers.

A test reference of Bush's `As We May Think' \cite{bush:1945:awmt}, to indicate how references are set up.

There is no need to save this file when editing as Overleaf auto-saves the document. Use the green `Recompile' button in the upper menu bar to show or refresh the PDF. The output PDF is not saved into the project but can be downloaded, for reading outside overleaf, using a button at the right-hand end of the Overleaf menu-bar. To download the project files, use the home button in the menu bar to exit the project. Download options are available in your project listing.

Depending on your type of Overleaf account you will be able to share access with fellow authors---see Overleaf's documentation for details. Shared access also allows for leaving comments in the project.

\subsection{Do Not Remove Boilerplate Code}
The TEX document includes code to generate sections such as the the copyright block. Do not remove these. Authors of accepted papers will receive sections of code to replace these and customise the final paper accordingly.

Pay attention to the code comments about author information and add/remove authors as necessary; there must be at least one author. It is desirable that all/most authors have an ORCID ID as this is replacing the need for explicit emails (that my change as researcher move around. If an ORCID is supplied that author's name is made the anchor text of a web link to their ORCID page.

Macro codes are offered (\texttt{e.g. authornote}) which may be used as well as indicating the corresponding author(s) for communication during publication.

\section{Real Introduction}

In this report, we will explore the system architecture, implementation and performance evaluation of the course project.

The goal of the course project is to “implement a basic load shed- ding strategy for a CEP system to handle bursty workloads while ensuring low-latency processing based on codebase OpenCEP” [1]. The OpenCEP codebase is an unlicensed github repository which belongs to the user ilya-kolchinsky [2]. However, it claims to be a library and a framework for implementing complex event processing although it is impossible to install because the project is missing a pyproject.toml and setup.py file. Nevertheless, we made a copy of the project and fixed all the bugs in it to get started.

The project has been implemented using Python and the uv astral package and project manager, which is written in Rust [3]. The project uses git for version control and vscode as the text editor. To run the project, you must first clone the repository with git and install all required python packages using the command uv sync. Then to run the pattern recognition on some citibike test data, you need to run:

\begin{lstlisting}
uv run cep-run \
    --patterns config/citibike_patterns.yaml \
    --input test_small.csv \
    --mode citibike --verbose
\end{lstlisting}

The goal of this project is to identify hot routes in CitiBike travel data, that is, identifying chains of bicycle trips that end at busy stations. This is a very important problem because CitiBike operators move hundreds of bicycles every day to balance the amount of bikes at each station. By identifying these hot paths, we can improve operational efficiency and maintain high availability of bikes.

\section{System Architecture}

The system architecture is built around dynamically adjusting internal parameters to handle resource exhaustion while maintaining low latency. The system reads the state of the CEP engine and makes load shedding decisions based on it.

The system separates event ingestion, pattern processing and system monitoring into separate components. Events are filtered before CEP processing according to a four-tier priority queue. The priorities from highest to lowest are: partial match protection, target station events, noise elimination, and adaptive sampling. System monitoring provides data based on event processing latency which determines parameter adjustment. The system is optimized to maintain low latency rather than high throughput. When latency exceeds targets, sampling rate is lowered to reduce latency. The architecture assumes regular resource constraints.


\section{Implementation}

\subsection{OpenCEP Framework Optimizations}

\textbf{Kleene Closure Performance.} The original OpenCEP implementation generates the full powerset ($2^n$ combinations) for Kleene closure operators, which caused an exponential state explosion in our case. With up to 100 partial matches in the hot paths pattern, execution would simply freeze. Our solution pre-filters matches by grouping attributes (bike ID first), then iterates from maximum to minimum size while collecting temporally valid sequences. The complexity drops from $O(2^n)$ to $O(n \cdot m)$ where $m$ is the Kleene size bound. The trade-off is that we sacrifice exhaustive enumeration for linear complexity, essentially prioritizing longest matches over generating all possible combinations.

\textbf{Time Window Correctness.} OpenCEP's Event class only tracked a single timestamp, which breaks time window calculations for events with duration like bike trips. We added an end timestamp getter to the DataFormatter API and modified Event to track both min\_timestamp (trip start) and max\_timestamp (trip end). This ensures that the 1-hour window evaluation works correctly for trip chains.

\subsection{Hot Paths Pattern}

The pattern detects bike trip chains ending at NYC's top-3 busiest stations in the August 2018 dataset (402, 435, 497). We use a sequence pattern with three main constraints: same bike ID across all trips in the chain, spatial chaining where each trip's end station matches the next trip's start station, and the final trip must end at one of our target stations. All of this happens within a 1-hour window. Setting the maximum size to 10 on the Kleene closure combined with our greedy matching approach keeps the state explosion under control while still catching all valid hot paths.

\subsection{Pattern-Aware Load Shedding}

The load shedding uses a four-tier priority hierarchy:
\begin{enumerate}
  \item \textbf{Partial Match Protection}: Events with bike ID or station ID in active partial matches are never dropped (queried from CEP tree every 0.5s), ensuring that patterns can complete.
  \item \textbf{Target Station Events}: Events to/from stations 402, 435, 497 are always kept, biasing toward pattern discovery.
  \item \textbf{Same-Station Elimination}: Round trips (start = end station) are dropped as noise.
  \item \textbf{Adaptive Sampling}: When latency ratio exceeds 1.5 or queue fill exceeds 70\%, remaining events are probabilistically dropped. The sampling rate adapts according to load: the aggressive mode (ratio $> 2.0$) reduces to 50\%, the moderate (ratio $> 1.5$) to 70\%, the relaxed (ratio $< 0.5$) increases toward 100\%.
\end{enumerate}

This ordering prioritizes recall (protecting active patterns) over throughput, trading latency for pattern completeness under load.

\subsection{Threading Architecture}

Three threads handle different aspects: (1) \textbf{Feeder} streams events from CSV, applies load shedding, and queues in a bounded queue (1000 events); (2) \textbf{CEP Engine} (main thread) consumes from the queue and performs pattern matching; (3) \textbf{Monitor} reports real-time stats without blocking CEP. The bounded queue provides backpressure when CEP cannot keep pace.

Threading is necessary to maintain state continuity across the entire stream. Naive chunked processing loses partial matches between chunks, which drastically reduces recall. Our architecture feeds one continuous stream to preserve all CEP state while enabling asynchronous I/O and monitoring.

\section{Performance Evaluation}

We evaluate recall, latency, throughput, and scalability using August 2018 CitiBike data (20,000 events sampled from 1.98M events over 31 days). To stress-test load shedding, we use an aggressive 5ms base latency instead of the standard 50ms, forcing the system to shed load even with moderate event rates. A baseline run with no load shedding establishes 100\% recall and system capacity. We then test five latency bounds (10\%, 30\%, 50\%, 70\%, 90\% of baseline latency) where the system adaptively sheds load to meet targets. Metrics tracked include throughput (events/sec), recall (\% of baseline matches), memory (MB peak), and drop rate (\%).

\subsection{Results}

% TODO: Run baseline + latency bounds and fill in table
Table~\ref{tab:results} shows performance across configurations.

\begin{table*}
  \caption{Performance Evaluation Results (20,000 events, August 2018 CitiBike data, 5ms aggressive base latency)}
  \label{tab:results}
  \begin{tabular}{lrrrrrr}
    \toprule
    Config   & Events & Matches & Recall & Drop   & Throughput & Memory \\
             & Proc.  & Found   & (\%)   & Rate   & (ev/s)     & (MB)   \\
    \midrule
    Baseline & 20000  & 182     & 100.0  & 0\%    & 148        & 41.7   \\
    10\% LB  & 11251  & 117     & 64.3   & 43.7\% & 325        & 41.7   \\
    30\% LB  & 11927  & 133     & 73.1   & 40.4\% & 321        & 41.7   \\
    50\% LB  & 12910  & 165     & 90.7   & 35.4\% & 271        & 41.7   \\
    70\% LB  & 13470  & 158     & 86.8   & 32.6\% & 257        & 41.8   \\
    90\% LB  & 14719  & 150     & 82.4   & 26.4\% & 241        & 41.8   \\
    \bottomrule
  \end{tabular}
\end{table*}

\textbf{Recall vs. Latency Trade-off.} Under aggressive latency constraints (5ms base, 10× stricter than standard 50ms), the system demonstrates recall-latency trade-offs. At the 10\% bound (0.5ms target), aggressive load shedding drops 43.7\% of events, achieving 64.3\% recall (117/182 matches). As latency bounds relax, recall improves: 30\% bound yields 73.1\% recall (133 matches, 40.4\% drop), 50\% bound reaches the highest recall at 90.7\% (165 matches, 35.4\% drop), while 70\% and 90\% bounds show 86.8\% and 82.4\% recall respectively. The non-monotonic pattern at higher bounds reflects the stochastic nature of adaptive sampling under load—when the system aggressively drops events to meet latency targets, some pattern chains are interrupted unpredictably. This validates that pattern-aware load shedding successfully protects critical events through partial match protection, though extreme latency constraints introduce variability in recall.

\textbf{Scalability.} The baseline processes 20,000 events in 135.0 seconds at 148 events/sec with 41.7 MB peak memory. Under aggressive load shedding (10\% bound), throughput increases to 325 events/sec (+120\%) by dropping 43.7\% of events, achieving 2.2× faster processing at the cost of 35.7\% recall loss. At relaxed bounds, throughput decreases to 241-321 events/sec as fewer events are dropped. Memory footprint remains constant at 41.7-41.8 MB across all configurations, demonstrating efficient state management. The system detects trip chains of up to 4 trips within the 1-hour temporal window, validating the temporal filtering and cleanup strategy. OpenCEP's sequential-only evaluation prevents multi-core parallelism, limiting CPU utilization to a single core.

\textbf{Optimization Impact.} The Kleene closure optimization ($O(2^n)$ $\rightarrow$ $O(n \cdot m)$) transformed an unusable system (which froze on powerset generation) into one processing 148 events/sec and detecting 182 patterns. This validates our complexity analysis and demonstrates why algorithmic optimization is critical for practical CEP applications.

\subsection{Limitations}

\textbf{Dataset Scale.} The August 2018 dataset contains approximately 1.98M events over 31 days. While the natural event rate is low, we artificially stress-tested the system using aggressive 5ms base latency (10× stricter than standard 50ms). This successfully triggered load shedding across all latency bounds, with drop rates ranging from 26.4\% to 43.7\%. However, this artificial constraint limits real-world applicability—production systems would likely use higher latency targets. Future work should evaluate with time-compressed replay at realistic latency bounds to better understand production performance.

\textbf{OpenCEP Constraints.} Sequential-only evaluation limits scalability. The lack of a public API for partial match state required fragile internal tree traversal. These constraints reduce portability and prevent multi-core scaling.

%%
%% The acknowledgments section is defined using the "acks" environment
%% (and NOT an unnumbered section). This ensures the proper
%% identification of the section in the article metadata, and the
%% consistent spelling of the heading.
\begin{acks}
  Acknowledgements go here. Delete enclosing begin/end markers if there are no acknowledgements.
\end{acks}

%%
%% The next two lines define the bibliography style to be used, and
%% the bibliography file.
\bibliographystyle{ACM-Reference-Format}
\bibliography{references.bib}

%%
\end{document}
\endinput
%%
%% End of file `main.tex'.
